\documentclass[11pt]{article}
\usepackage[a4paper,margin=1in]{geometry}
\usepackage{amsmath}
\usepackage{booktabs}

\title{Why the Micro Simulator Appears Faster Than MicroMacro}
\author{}
\date{}

\begin{document}
\maketitle

\section*{Scope}
This note is a theoretical explanation for a synthetic two-layer stochastic network model used for software testing and visualization. It explains why the \texttt{Micro} simulator can show faster spread dynamics than \texttt{MicroMacro} under the current setup.

\section*{Observed Setup (from your config)}
\begin{itemize}
\item Communities: $k=4$
\item Community size: $N=100$
\item Macro topology: chain
\item Inter-community links: $W_{ij}=1$ for adjacent communities
\item Intra-community graph: complete
\item $\beta=0.005,\ \gamma=0,\ \tau_{\text{micro}}=1,\ T=1$
\end{itemize}

\section*{Core Reason 1: Different Inter-Community Rate Laws}
\textbf{Micro (edge-level SSA):} cross-community infection is driven by actual infected--susceptible inter-community edges in the full graph. Once a bridge endpoint is infected, transmission over that active edge occurs at rate approximately $\beta$ (per active edge).

\textbf{MicroMacro (community-level coupling):} macro transfer hazard is
\[
\lambda_{ij} \;=\; \beta_{\text{macro}}\,T\,W_{ij}\,\frac{I_i}{N_i}\,\frac{S_j}{N_j}.
\]
Early in an outbreak, $I_i/N_i$ is small. With $N_i=100$ and initially $I_i=1$, the prefactor is $0.01$, so macro transfer hazards are strongly suppressed at early times.

Therefore, \texttt{MicroMacro} requires source-community prevalence to build before exporting effectively, while \texttt{Micro} can export as soon as bridge-node infection appears.

\section*{Core Reason 2: Coupling via Time-Split Integration Adds Lag}
\texttt{MicroMacro} advances with interval splitting (\(\tau_{\text{micro}}=1\)) and uses midpoint hazard accumulation for macro-event timing. This introduces a finite-step approximation error and a response delay when hazards increase rapidly inside an interval.

In contrast, \texttt{Micro} is event-driven on the full graph without macro/micro splitting, so inter-community events occur at exact SSA event times.

\section*{Core Reason 3: Chain Topology Amplifies Delays}
With a chain macro graph and one inter-community link between neighbors, downstream communities can only activate through sequential transfer:
\[
0 \rightarrow 1 \rightarrow 2 \rightarrow 3.
\]
Any lag in one hop propagates to later communities. This stacking effect is stronger in \texttt{MicroMacro} because each hop depends on prevalence-weighted macro hazards and threshold crossing.

\section*{Core Reason 4: Complete Intra-Community Graph Speeds Bridge Activation in Micro}
Inside each community, the complete graph causes very fast mixing. In \texttt{Micro}, this quickly infects bridge-relevant nodes, enabling earlier direct export over real inter-community edges.

In \texttt{MicroMacro}, fast intra-community growth does not immediately translate to export; export remains filtered through the normalized macro hazard.

\section*{Secondary Visual Effect}
There can also be a plotting-level right-shift if one output stream starts recording at $t=\tau_{\text{micro}}$ while the other includes $t=0$. That does not fully explain the dynamic gap, but it can make \texttt{MicroMacro} appear additionally delayed in plots.

\section*{Summary}
The faster dynamics in \texttt{Micro} are expected from model structure, not just randomness:
\begin{itemize}
\item edge-level immediate cross-community transmission vs. prevalence-normalized macro transfer;
\item exact full-graph event timing vs. split-step macro coupling;
\item chain-hop delay accumulation under sparse inter-community connectivity.
\end{itemize}
Under your current parameters (especially \(N=100\), \(W_{ij}=1\), chain coupling), these effects systematically bias \texttt{MicroMacro} toward slower apparent spread.

\end{document}
